% Pursuing Stacks (cotes 134-2 to 134-8): archive cote 134-3, site page 58 (PDF page 59), typescript page 58.
% typescript. Draft compared: ps2.tex:2932-3005.
% AI-transcribed by Claude Opus 5.5 (claude-opus-5-5, adaptive thinking, effort high), Message Batches API, batch ps1 (msgbatch_01Rv5hCaNTtUt4JHx1yqoHPL), 2026-11-12; not yet checked by eye.
\scanpage{134-3/58}
\mspage{58}

\margin{*} tute in A\^{}. We could formalize it as follows:\hand{..} \hand{(42).}

\margin{\unsure{cor. inst ci-dessous}} \deleted{\uline{Corollary:} In the proposition above, the condition that $\hand{\alpha}(\hand{\Delta}_1)$ be aspheric}
\deleted{over $e_{A\hat{}}$ is a consequence of the following assumptions:}
\deleted{a) The condition T 3 is valid in A , with the stronger assumption \xout{th}}
\deleted{that I is aspheric over the final object e of A\^{}, and moreover :}
\deleted{b) for every $a \hand{\in} \mathrm{Ob}\, A$ and $u : a \hand{\longrightarrow} I$, let $a_u = u^{-1}(e_o)$ (subobject}
\deleted{of a in A\^{}), $\hand{\Delta}_{a_u} = \varinjlim_{b \text{ in } A_{/a_u}} \hand{\Delta}_b$ , and $\hand{\Delta}_{a,u} = \mathrm{Im}\,(\hand{\Delta}_{a_u} \hand{\to} \hand{\Delta}_a)$.}

\deleted{We assume that Ob $\hand{\Delta}_{a,u}$ is a \uline{crible} in $\hand{\Delta}_a$ , say $C_u$ .}

\deleted{Indeed, interpreting $\hand{\alpha}(\hand{\Delta}_1)$ as the functor $a \hand{\mapsto} \mathrm{Crib}\, \hand{\Delta}_a$ , we}
\deleted{define a homotopy}
\deleted{\[ h : I \times F\deleted{\unsure{(\alpha(\Delta_1))}} \hand{\to} F \deleted{\unsure{\alpha(\Delta_1)}} \quad \text{(where } F = \hand{\alpha}(\hand{\Delta}_1)\text{ )} \]}
\deleted{from the identity $\mathrm{id}_F$ to the constant map $c : F \hand{\to} F$ associating to every}
\deleted{crible in some $\hand{\Delta}_a$ the empty crible in the same, by taking for every a}
\deleted{in A the map}
\deleted{\[ (u, C_o) \hand{\longmapsto} C_u \hand{\cap} C_o \quad . \]}
\deleted{It follows that $F \hand{\unsure{\longmapsto|\to}} e$ is a weak equivalence, and the same argument in}
\deleted{any induced category $A_{/a}$ shows that this is universally so i.e. $F \hand{\to} e$ is}
\deleted{aspheric , O.K.}

\deleted{In the case \inserted{above,} $A = \hand{\tilde{\Delta}}$, we get indeed that for $u : a \hand{\to} I$, $a_u$ is \inserted{\unsure{empty or,}} in A}
\deleted{and $\hand{\Delta}_{a_u} \hand{\to} \hand{\Delta}_a$ is a full embedding, turning $\hand{\Delta}_{a_u}$ into a crible in $\hand{\Delta}_a$ ,}
\deleted{thus \xout{a) and} b) is satisfied (with the usual choice $I = \hand{\tilde{\Delta}}_1$ of course).}

21.3.

37. By the end of the notes two days ago, there was the pretty strong
impression of repeating the same argument all over again, in very similar
situations. When I tried to pin down this feeling, first thing that oc-%
cured to me was that the \inserted{homotopy} equivalence I was after $A_{/C} \hand{\longrightarrow} C$ (in the propo-%
sition of \over{\unsure{p}}{\S}.\over{\unsure{5}}{3}6), when concerned with a\xout{n} general functor $A \hand{\longrightarrow} (\mathrm{Cat})$,
and the criterion obtained, was applicable to the situation I was concer-%
ned with at the very start when defining test categories, namely when
looking at the \uline{canonical} functor $A \hand{\longrightarrow} (\mathrm{Cat})$ given by $a \hand{\longmapsto} A_{/a}$ ; and \inserted{that this} gave
a criterion in this case for A to be a test-category in the wider sense,
which I had overlooked when peeling out the caracterisation of test cate-%
gories (cf. theorem of p. 46). Finally it bec\over{o}{a}me\unsure{s} clear that it \over{was}{is} about
% NOTE: the right edge of the next line is cut; the typed word ends "repet" at the edge
time to recast from scratch the asphericity story, and (by one more repet%
tition) tell it anew in a way stripped from it's repetitive features\over{.}{!}
